\chapter*{Resumen}
\addcontentsline{toc}{chapter}{Resumen}

% TODO: Escribir resumen en español (~250 palabras)
% Puntos a cubrir:
%   - Pregunta de investigación: ¿puede el ML + NLP predecir mejor el EMBI de Ecuador que métodos econométricos?
%   - Datos: 3.047 días de trading, 2013-2025, fuentes GDELT v2, Bloomberg/Yahoo Finance, macroeconomía local
%   - Métodos: ARIMA baseline, XGBoost con ablación, Granger, Chow/CUSUM, Autoencoder NLP
%   - Resultado principal: Grupo C (AR + NLP) alcanza RMSE 49.01 pb vs 107.48 pb del ARIMA (reducción 54.4%)
%   - Hallazgo exogeneidad: variables locales no mejoran el modelo; en crisis el desempleo local irrumpe
%   - Hallazgo NLP: la cobertura mediática (volumen) aporta señal; el sentimiento no
%   - Hallazgo autoencoder: la representación NLP es esencialmente lineal (PCA suficiente)

\vspace{1cm}

\noindent \textbf{Palabras clave:} riesgo país, EMBI, Ecuador, XGBoost, GDELT, NLP, quiebres estructurales, análisis de regímenes.

\vspace{1cm}

\chapter*{Abstract}
\addcontentsline{toc}{chapter}{Abstract}

% TODO: Traducir resumen al inglés

\vspace{1cm}

\noindent \textbf{Keywords:} country risk, EMBI, Ecuador, XGBoost, GDELT, NLP, structural breaks, regime analysis.
